% 局部参考模式，不是主模板。仅选择所需 frame 并替换全部占位词。
% 不要直接 include 本文件，不要复制进主文件的导言区。

% ---------- 重要逻辑节点：延续现有整页大字风格 ----------
\begin{frame}
    \centering
    {\LARGE\bfseries 在此提出本篇论文真正需要回答的关键问题\par}
    \vfill
    {\small 用一句话提示连接当前障碍与下一步推导的工具\par}
\end{frame}

% ---------- 正式结果：先说明前提，再解释结论 ----------
% Source: 此处记录原文 theorem/equation/page 定位。
\begin{frame}{本页要建立的核心结论}
    \small
    \begin{block}{设定与条件}
        在此给出本结论涉及的对象、参数、定义域和必要假设。
    \end{block}
    \begin{alertblock}{结果（原文编号）}
        在此写出经过核对的完整数学陈述。
    \end{alertblock}
    \textbf{含义：}说明结论为何解决当前问题，以及下一步将如何使用它。
\end{frame}

% ---------- 连续推导页：按一个变换拆页 ----------
\begin{frame}{推导的当前步骤及目标}
    \small
    \textbf{起点：}用上一页得到的关系或明确前提开始。
    \begin{block}{本步推导}
        用 aligned/align* 写真实推导，非显然步骤标明依据。
    \end{block}
    \textbf{本步依据：}说明引理、换元、最优化或积分交换的适用条件。

    \vspace{0.5ex}
    \alert{用一句话说明本步结论怎样进入后续论证。}
\end{frame}

% ---------- 图片与解释：使用 main.tex 已有 graphicspath ----------
% Source: 原文 Figure 编号、PDF 页码、原文件和裁切方式。
\begin{frame}{图片检验的具体问题}
    \small
    \begin{columns}[T,onlytextwidth]
        \begin{column}{0.55\textwidth}
            \centering
            \includegraphics[width=\linewidth,height=0.57\textheight,keepaspectratio]
                {paperX/descriptive-name.pdf}

            {\footnotesize 原文 Figure 编号及必要的设置说明\par}
        \end{column}
        \begin{column}{0.42\textwidth}
            \begin{block}{如何读图}
                说明坐标、指标方向、比较条件及关键观测。
            \end{block}
            \alert{只陈述图片和原文实际支持的结论。}
        \end{column}
    \end{columns}
\end{frame}
